\chapter{PyTorch入门}

\intro{PyTorch 是目前最流行的深度学习框架之一，以其动态计算图、Pythonic 的编程风格和强大的 GPU 加速能力而受到研究者广泛喜爱。本章从零开始介绍 PyTorch 的核心概念：张量操作、自动求导、神经网络模块、优化器、数据加载，以及完整的训练流程。}

\section{张量基础}

\subsection{创建张量}

张量（Tensor）是 PyTorch 中的核心数据结构，类似于 NumPy 的 ndarray，但额外支持 GPU 加速和自动微分。

\begin{mycode}{张量创建与基本操作}
\begin{lstlisting}[language=Python]
import torch

# 创建张量
a = torch.tensor([1, 2, 3])
b = torch.zeros(2, 3)
c = torch.ones(3, 4)
d = torch.randn(3, 3)      # 标准正态分布
e = torch.arange(0, 10, 2) # 0到10，步长2
f = torch.eye(3)           # 3x3单位矩阵

print(f"a: {a}, 形状: {a.shape}")
print(f"b:\n{b}")
print(f"d:\n{d}")

# 从NumPy转换
import numpy as np
np_arr = np.array([[1, 2], [3, 4]])
torch_arr = torch.from_numpy(np_arr)
back_to_np = torch_arr.numpy()
\end{lstlisting}
\end{mycode}

\subsection{张量属性与运算}

\begin{mycode}{张量属性与运算}
\begin{lstlisting}[language=Python]
import torch

x = torch.randn(2, 3, 4)
print(f"形状: {x.shape}")
print(f"维度数: {x.ndim}")
print(f"元素总数: {x.numel()}")
print(f"数据类型: {x.dtype}")
print(f"所在设备: {x.device}")

# 基本运算
a = torch.tensor([1., 2., 3.])
b = torch.tensor([4., 5., 6.])
print(f"加法: {a + b}")
print(f"乘法(逐元素): {a * b}")
print(f"点积: {torch.dot(a, b)}")

# 矩阵运算
A = torch.randn(3, 4)
B = torch.randn(4, 2)
C = A @ B  # 矩阵乘法
print(f"矩阵乘法结果形状: {C.shape}")

# 维度操作
x = torch.randn(2, 3, 4)
print(f"reshape: {x.reshape(6, 4).shape}")
print(f"转置: {x.transpose(0, 2).shape}")
print(f"permute: {x.permute(2, 1, 0).shape}")
print(f"unsqueeze: {x.unsqueeze(0).shape}")  # 增加维度
print(f"squeeze: {x.unsqueeze(0).squeeze(0).shape}")

# 拼接
a = torch.randn(2, 3); b = torch.randn(2, 3)
print(f"cat(dim=0): {torch.cat([a, b], dim=0).shape}")
print(f"cat(dim=1): {torch.cat([a, b], dim=1).shape}")
print(f"stack: {torch.stack([a, b], dim=0).shape}")
\end{lstlisting}
\end{mycode}

\subsection{索引与切片}

张量的索引和切片与 NumPy 类似，支持高级索引。关键概念：\textbf{视图（view）与副本（copy）}——大多数切片操作返回视图（共享底层数据），使用 \texttt{.clone()} 创建独立副本。

\begin{mycode}{张量索引与切片}
\begin{lstlisting}[language=Python]
import torch

x = torch.arange(24).reshape(2, 3, 4)
print(f"张量形状: {x.shape}")

# 基本索引
print(f"x[0]:\n{x[0]}")           # 第一个3x4矩阵
print(f"x[0, 1]: {x[0, 1]}")     # 第一矩阵的第二行
print(f"x[:, :2, :]:\n{x[:, :2, :]}")  # 所有批次的前两行

# 花式索引
indices = torch.tensor([0, 2])
print(f"x[:, :, [0,2]]:\n{x[:, :, indices]}")

# 布尔索引
mask = x > 10
print(f"大于10的元素: {x[mask]}")

# view与副本
y = x.view(6, 4)  # 视图（共享数据）
z = x.clone()      # 独立副本
y[0, 0] = 999
print(f"修改view后原始张量x[0,0,0]: {x[0,0,0]}")  # 已被修改!
\end{lstlisting}
\end{mycode}

\section{GPU加速}

PyTorch 的一大优势是无缝的 GPU 支持。张量可以在 CPU 和 GPU 之间轻松移动：

\begin{mycode}{GPU操作}
\begin{lstlisting}[language=Python]
import torch

# 检查GPU可用性
print(f"CUDA可用: {torch.cuda.is_available()}")
if torch.cuda.is_available():
    print(f"GPU数量: {torch.cuda.device_count()}")
    print(f"当前GPU: {torch.cuda.get_device_name(0)}")

# 在指定设备上创建张量
device = torch.device('cuda' if torch.cuda.is_available() else 'cpu')

x_cpu = torch.randn(1000, 1000)
x_gpu = x_cpu.to(device)       # 移动到GPU
x_back = x_gpu.cpu()           # 移回CPU

# 直接在目标设备创建
y = torch.randn(1000, 1000, device=device)

# GPU加速矩阵乘法
import time
A = torch.randn(5000, 5000, device=device)
B = torch.randn(5000, 5000, device=device)

if device.type == 'cuda':
    torch.cuda.synchronize()
t0 = time.time()
C = A @ B
if device.type == 'cuda':
    torch.cuda.synchronize()
print(f"GPU 矩阵乘法时间: {time.time()-t0:.3f}s")

# 内存管理
print(f"GPU内存分配: {torch.cuda.memory_allocated()/1e9:.2f} GB")
print(f"GPU内存缓存: {torch.cuda.memory_reserved()/1e9:.2f} GB")
torch.cuda.empty_cache()  # 清理缓存
\end{lstlisting}
\end{mycode}

\section{自动求导}

自动求导（Autograd）是 PyTorch 的核心功能，它自动计算计算图中任意参数的梯度。

\subsection{计算图}

PyTorch 在每次前向传播时动态构建计算图。设置 \texttt{requires\_grad=True} 的张量会记录其上的所有操作，调用 \texttt{.backward()} 后自动计算梯度。

\begin{mycode}{自动求导基础}
\begin{lstlisting}[language=Python]
import torch

# 需要梯度的张量
x = torch.tensor([2.0, 3.0], requires_grad=True)
w = torch.tensor([1.0, -1.0], requires_grad=True)
b = torch.tensor([0.5], requires_grad=True)

# 前向计算: y = w·x + b
y = torch.dot(w, x) + b

# 反向传播
y.backward()

print(f"y = {y.item()}")
print(f"dy/dw = {w.grad}")  # 应为 x = [2, 3]
print(f"dy/dx = {x.grad}")  # 应为 w = [1, -1]
print(f"dy/db = {b.grad}")  # 应为 1

# 多次反向传播需要 retain_graph=True
x = torch.tensor(3.0, requires_grad=True)
y = x ** 2
y.backward(retain_graph=True)
print(f"dy/dx = {x.grad}")  # 2*x = 6

# 高阶导数
x = torch.tensor(2.0, requires_grad=True)
y = x ** 3
grad1 = torch.autograd.grad(y, x, create_graph=True)[0]
print(f"一阶导数 (3x^2 at x=2): {grad1.item()}")  # 12
grad2 = torch.autograd.grad(grad1, x)[0]
print(f"二阶导数 (6x at x=2): {grad2.item()}")    # 12
\end{lstlisting}
\end{mycode}

\subsection{梯度在神经网络中的应用}

\begin{mycode}{使用autograd训练简单模型}
\begin{lstlisting}[language=Python]
import torch

torch.manual_seed(42)
m, d = 200, 5
X = torch.randn(m, d)
true_w = torch.tensor([2.0, -1.0, 3.0, -0.5, 1.5])
y = X @ true_w + 2.0 + 0.3 * torch.randn(m)

w = torch.randn(d, requires_grad=True)
b = torch.zeros(1, requires_grad=True)
lr = 0.05

for epoch in range(500):
    # 前向
    y_pred = X @ w + b
    loss = 0.5 * torch.mean((y_pred - y) ** 2)

    # 反向
    loss.backward()

    # 参数更新（手动，之后会介绍优化器）
    with torch.no_grad():
        w -= lr * w.grad
        b -= lr * b.grad
        w.grad.zero_()
        b.grad.zero_()

    if epoch % 100 == 0:
        print(f"Epoch {epoch:3d}, Loss: {loss.item():.6f}")

print(f"\n真实w: {true_w}")
print(f"学习w: {w.data}")
print(f"学习b: {b.item():.4f}")
\end{lstlisting}
\end{mycode}

\begin{figure}[H]
\centering
\begin{tikzpicture}[
    mynode/.style={rectangle, draw=blue!60, fill=blue!5, minimum width=1.5cm, minimum height=0.7cm, rounded corners=2pt},
    mygrad/.style={rectangle, draw=red!60, fill=red!5, minimum width=1.5cm, minimum height=0.7cm, rounded corners=2pt},
    myedge/.style={->, >=stealth, thick, blue!70},
    mygrada/.style={->, >=stealth, thick, red!70, bend right},
    mylabel/.style={font=\footnotesize}
]
    % Input
    \node[mynode] (X) at (0, 3) {$\mathbf{X}$};
    \node[mynode] (W) at (0, 1.5) {$\mathbf{w}$};
    \node[mynode] (b) at (0, 0) {$b$};

    % Operations
    \node[mynode, draw=green!60, fill=green!5] (mul) at (3, 2) {$\cdot$};
    \node[mynode, draw=green!60, fill=green!5] (add) at (5.5, 1.5) {$+$};
    \node[mynode, draw=green!60, fill=green!5] (mse) at (8, 1.5) {MSE};

    % Output
    \node[mygrad] (L) at (10, 1.5) {Loss};

    % Forward edges
    \draw[myedge] (X) |- (mul);
    \draw[myedge] (W) -- (mul);
    \draw[myedge] (mul) -- (add);
    \draw[myedge] (b) |- (add);
    \draw[myedge] (add) -- (mse);
    \draw[myedge] (mse) -- (L);

    % Gradient edges
    \draw[mygrada] (L) -- (mse);
    \draw[mygrada] (mse) -- (add);
    \draw[mygrada] (add) -- (mul);
    \draw[mygrada] (add) |- (b);
    \draw[mygrada] (mul) -- (W);
    \draw[mygrada] (mul) |- (X);

    % Labels
    \node[mylabel, red!70] at (4, 0.2) {反向传播（梯度流）};
    \node[mylabel, blue!70] at (6.5, 3.5) {前向传播（计算流）};
\end{tikzpicture}
\caption{PyTorch 自动求导的计算图。蓝色实线是前向传播，红色虚线是反向传播的梯度流。每个节点记录局部导数，通过链式法则自动合成完整梯度。}
\label{fig:computation-graph}
\end{figure}

\section{nn.Module：构建神经网络}

\texttt{torch.nn.Module} 是 PyTorch 中构建神经网络的基类。它封装了参数、前向传播逻辑以及参数管理。

\begin{mycode}{定义神经网络模块}
\begin{lstlisting}[language=Python]
import torch
import torch.nn as nn

class SimpleNet(nn.Module):
    def __init__(self, input_dim, hidden_dim, output_dim):
        super().__init__()
        self.fc1 = nn.Linear(input_dim, hidden_dim)
        self.relu = nn.ReLU()
        self.fc2 = nn.Linear(hidden_dim, hidden_dim)
        self.fc3 = nn.Linear(hidden_dim, output_dim)

    def forward(self, x):
        x = self.fc1(x)
        x = self.relu(x)
        x = self.fc2(x)
        x = self.relu(x)
        x = self.fc3(x)
        return x

# 创建模型
model = SimpleNet(input_dim=10, hidden_dim=32, output_dim=3)
print(model)

# 查看参数
for name, param in model.named_parameters():
    print(f"{name}: {param.shape}")

# 总参数量
total_params = sum(p.numel() for p in model.parameters())
trainable_params = sum(p.numel() for p in model.parameters()
                       if p.requires_grad)
print(f"\n总参数: {total_params}, 可训练: {trainable_params}")
\end{lstlisting}
\end{mycode}

\subsection{常用层}

\begin{mycode}{常用神经网络层}
\begin{lstlisting}[language=Python]
import torch.nn as nn

# 线性层（全连接层）
fc = nn.Linear(in_features=128, out_features=64)

# 激活函数
relu = nn.ReLU()
gelu = nn.GELU()
sigmoid = nn.Sigmoid()
tanh = nn.Tanh()
leaky_relu = nn.LeakyReLU(negative_slope=0.01)

# 卷积层
conv1d = nn.Conv1d(in_channels=3, out_channels=16, kernel_size=3)
conv2d = nn.Conv2d(in_channels=3, out_channels=32,
                   kernel_size=3, padding=1, stride=2)

# 池化层
maxpool = nn.MaxPool2d(kernel_size=2)
avgpool = nn.AvgPool2d(kernel_size=2)
adaptive_pool = nn.AdaptiveAvgPool2d((1, 1))

# 归一化层
batch_norm = nn.BatchNorm1d(64)
layer_norm = nn.LayerNorm(128)

# 正则化
dropout = nn.Dropout(p=0.5)

# 循环层
lstm = nn.LSTM(input_size=64, hidden_size=128,
               num_layers=2, batch_first=True)
gru = nn.GRU(input_size=64, hidden_size=128, batch_first=True)

# 嵌入层
embedding = nn.Embedding(num_embeddings=10000, embedding_dim=256)

# 打印参数量
print(f"Conv2d参数量: {sum(p.numel() for p in conv2d.parameters())}")
print(f"LSTM参数量: {sum(p.numel() for p in lstm.parameters())}")
\end{lstlisting}
\end{mycode}

\subsection{Sequential 容器}

对于简单的顺序模型，可以使用 \texttt{nn.Sequential} 简化代码：

\begin{mycode}{Sequential容器}
\begin{lstlisting}[language=Python]
import torch.nn as nn

# 方式1: 按顺序传入层
model = nn.Sequential(
    nn.Linear(28*28, 512),
    nn.ReLU(),
    nn.Dropout(0.2),
    nn.Linear(512, 256),
    nn.ReLU(),
    nn.Dropout(0.2),
    nn.Linear(256, 10)
)

# 方式2: 使用OrderedDict命名各层
from collections import OrderedDict
model_named = nn.Sequential(OrderedDict([
    ('fc1', nn.Linear(28*28, 512)),
    ('relu1', nn.ReLU()),
    ('fc2', nn.Linear(512, 256)),
    ('relu2', nn.ReLU()),
    ('fc3', nn.Linear(256, 10))
]))

print(model)
x = torch.randn(4, 28*28)  # batch of 4
out = model(x)
print(f"输出形状: {out.shape}")  # (4, 10)
\end{lstlisting}
\end{mycode}

\section{损失函数与优化器}

PyTorch 提供了丰富的损失函数和优化器：

\begin{mycode}{损失函数}
\begin{lstlisting}[language=Python]
import torch
import torch.nn as nn

# 回归损失
mse_loss = nn.MSELoss()
l1_loss = nn.L1Loss()
huber_loss = nn.SmoothL1Loss()

# 二分类损失
bce_loss = nn.BCELoss()              # 需要sigmoid后的输出
bce_logits = nn.BCEWithLogitsLoss()  # 包含sigmoid（更稳定）

# 多分类损失
ce_loss = nn.CrossEntropyLoss()  # 包含log_softmax

# 测试
# 回归示例
y_pred = torch.randn(32, 1)
y_true = torch.randn(32, 1)
print(f"MSE: {mse_loss(y_pred, y_true):.4f}")

# 分类示例
logits = torch.randn(32, 10)  # 原始分数（未归一化）
labels = torch.randint(0, 10, (32,))
print(f"交叉熵: {ce_loss(logits, labels):.4f}")
\end{lstlisting}
\end{mycode}

\begin{mycode}{优化器}
\begin{lstlisting}[language=Python]
import torch
import torch.nn as nn
import torch.optim as optim

model = nn.Linear(10, 3)

# SGD
sgd = optim.SGD(model.parameters(), lr=0.01, momentum=0.9,
                weight_decay=1e-4, nesterov=True)

# Adam
adam = optim.Adam(model.parameters(), lr=0.001,
                  betas=(0.9, 0.999), weight_decay=1e-4)

# AdamW (Adam + 解耦权重衰减，推荐使用)
adamw = optim.AdamW(model.parameters(), lr=0.001, weight_decay=0.01)

# RMSprop
rmsprop = optim.RMSprop(model.parameters(), lr=0.001, alpha=0.99)

# 学习率调度
optimizer = optim.AdamW(model.parameters(), lr=0.01)

# 阶梯衰减
scheduler_step = optim.lr_scheduler.StepLR(
    optimizer, step_size=30, gamma=0.1)

# 余弦退火
scheduler_cos = optim.lr_scheduler.CosineAnnealingLR(
    optimizer, T_max=100, eta_min=1e-6)

# OneCycleLR
scheduler_1cycle = optim.lr_scheduler.OneCycleLR(
    optimizer, max_lr=0.01, total_steps=1000)

# 训练中使用
for epoch in range(10):
    # ... 训练代码 ...
    optimizer.step()
    scheduler_cos.step()  # 每个epoch更新
    print(f"Epoch {epoch}: lr = {scheduler_cos.get_last_lr()[0]:.6f}")
\end{lstlisting}
\end{mycode}

\section{数据加载}

PyTorch 的 \texttt{torch.utils.data} 模块提供了标准的数据加载方式：

\begin{mycode}{Dataset与DataLoader}
\begin{lstlisting}[language=Python]
import torch
from torch.utils.data import Dataset, DataLoader, random_split
import numpy as np

# 自定义Dataset
class MyDataset(Dataset):
    def __init__(self, n_samples=1000, n_features=10):
        self.X = torch.randn(n_samples, n_features)
        self.y = (self.X.sum(dim=1) > 0).long()

    def __len__(self):
        return len(self.X)

    def __getitem__(self, idx):
        return self.X[idx], self.y[idx]

# 创建数据集
dataset = MyDataset(n_samples=1000)

# 划分训练/验证集
train_size = int(0.8 * len(dataset))
val_size = len(dataset) - train_size
train_dataset, val_dataset = random_split(
    dataset, [train_size, val_size])

# DataLoader
train_loader = DataLoader(train_dataset, batch_size=32,
                          shuffle=True, num_workers=0,
                          drop_last=True)
val_loader = DataLoader(val_dataset, batch_size=32,
                        shuffle=False, num_workers=0)

# 使用
for batch_idx, (data, target) in enumerate(train_loader):
    print(f"Batch {batch_idx}: data {data.shape}, target {target.shape}")
    if batch_idx >= 2:
        break
\end{lstlisting}
\end{mycode}

\subsection{内置数据集}

PyTorch 的 \texttt{torchvision} 提供了常用图像数据集：

\begin{mycode}{使用torchvision加载MNIST}
\begin{lstlisting}[language=Python]
import torch
from torch.utils.data import DataLoader
from torchvision import datasets, transforms

# 数据预处理
transform = transforms.Compose([
    transforms.ToTensor(),  # 转为张量并归一化到[0,1]
    transforms.Normalize((0.1307,), (0.3081,))  # MNIST均值/标准差
])

# 加载数据集
train_dataset = datasets.MNIST(
    root='./data', train=True, download=True,
    transform=transform)
test_dataset = datasets.MNIST(
    root='./data', train=False, download=True,
    transform=transform)

train_loader = DataLoader(train_dataset, batch_size=64,
                          shuffle=True)
test_loader = DataLoader(test_dataset, batch_size=1000,
                         shuffle=False)

# 查看数据
images, labels = next(iter(train_loader))
print(f"图像形状: {images.shape}")  # (64, 1, 28, 28)
print(f"标签形状: {labels.shape}")  # (64,)
print(f"标签值范围: [{labels.min()}, {labels.max()}]")
\end{lstlisting}
\end{mycode}

\section{训练循环}

一个完整的训练循环需要包含：训练模式设置、梯度清零、前向传播、损失计算、反向传播、参数更新和评估。

\begin{mycode}{完整的训练流程}
\begin{lstlisting}[language=Python]
import torch
import torch.nn as nn
import torch.optim as optim
from torch.utils.data import DataLoader, random_split

# 准备数据
class ToyDataset(torch.utils.data.Dataset):
    def __init__(self, n=3000, d=20, n_classes=5):
        self.X = torch.randn(n, d)
        self.y = torch.randint(0, n_classes, (n,))

    def __len__(self): return len(self.X)
    def __getitem__(self, idx): return self.X[idx], self.y[idx]

dataset = ToyDataset()
train_ds, val_ds = random_split(dataset, [2400, 600])
train_loader = DataLoader(train_ds, batch_size=64, shuffle=True)
val_loader = DataLoader(val_ds, batch_size=64, shuffle=False)

# 创建模型
class MLP(nn.Module):
    def __init__(self, in_dim, hidden, n_classes):
        super().__init__()
        self.net = nn.Sequential(
            nn.Linear(in_dim, hidden),
            nn.ReLU(),
            nn.Dropout(0.3),
            nn.Linear(hidden, hidden // 2),
            nn.ReLU(),
            nn.Dropout(0.3),
            nn.Linear(hidden // 2, n_classes)
        )

    def forward(self, x):
        return self.net(x)

device = torch.device('cuda' if torch.cuda.is_available() else 'cpu')
model = MLP(in_dim=20, hidden=128, n_classes=5).to(device)
criterion = nn.CrossEntropyLoss()
optimizer = optim.AdamW(model.parameters(), lr=0.001, weight_decay=1e-4)
scheduler = optim.lr_scheduler.CosineAnnealingLR(optimizer, T_max=50)

# 训练循环
def train_epoch(model, loader, criterion, optimizer, device):
    model.train()
    total_loss, correct, total = 0.0, 0, 0

    for data, target in loader:
        data, target = data.to(device), target.to(device)

        optimizer.zero_grad()
        output = model(data)
        loss = criterion(output, target)
        loss.backward()
        optimizer.step()

        total_loss += loss.item() * data.size(0)
        pred = output.argmax(dim=1)
        correct += pred.eq(target).sum().item()
        total += data.size(0)

    return total_loss / total, correct / total

@torch.no_grad()
def evaluate(model, loader, criterion, device):
    model.eval()
    total_loss, correct, total = 0.0, 0, 0

    for data, target in loader:
        data, target = data.to(device), target.to(device)
        output = model(data)
        loss = criterion(output, target)

        total_loss += loss.item() * data.size(0)
        pred = output.argmax(dim=1)
        correct += pred.eq(target).sum().item()
        total += data.size(0)

    return total_loss / total, correct / total

# 运行训练
epochs = 50
best_acc = 0.0
history = {'train_loss': [], 'train_acc': [], 'val_loss': [], 'val_acc': []}

for epoch in range(epochs):
    train_loss, train_acc = train_epoch(
        model, train_loader, criterion, optimizer, device)
    val_loss, val_acc = evaluate(
        model, val_loader, criterion, device)
    scheduler.step()

    history['train_loss'].append(train_loss)
    history['train_acc'].append(train_acc)
    history['val_loss'].append(val_loss)
    history['val_acc'].append(val_acc)

    if val_acc > best_acc:
        best_acc = val_acc
        torch.save(model.state_dict(), 'best_model.pt')

    if epoch % 10 == 0 or epoch == epochs - 1:
        print(f"Epoch {epoch:3d}: "
              f"Train Loss {train_loss:.4f}, Acc {train_acc:.3f} | "
              f"Val Loss {val_loss:.4f}, Acc {val_acc:.3f}")

print(f"\n最佳验证准确率: {best_acc:.3f}")

# 加载最佳模型
model.load_state_dict(torch.load('best_model.pt'))
\end{lstlisting}
\end{mycode}

\begin{figure}[H]
\centering
\begin{tikzpicture}[
    mybox/.style={rectangle, draw=blue!60, fill=blue!5, minimum width=3.5cm, minimum height=0.8cm, font=\small, rounded corners=3pt},
    myarrow/.style={->, >=stealth, thick, blue!70},
    mylabel/.style={font=\footnotesize}
]
    % Training loop flow
    \node[mybox] (data) at (0, 5) {加载数据};
    \node[mybox] (fwd) at (0, 3.2) {前向传播};
    \node[mybox] (loss) at (0, 1.4) {计算损失};
    \node[mybox] (zero) at (0, -0.4) {梯度清零};
    \node[mybox] (bwd) at (0, -2.2) {反向传播};
    \node[mybox] (step) at (0, -4.0) {参数更新};

    \node[mybox, draw=green!60, fill=green!5] (eval) at (4, 1) {验证/评估};
    \node[mybox, draw=red!60, fill=red!5] at (4, -0.5) {(每N个epoch)};

    \draw[myarrow] (data) -- (fwd);
    \draw[myarrow] (fwd) -- (loss);
    \draw[myarrow] (loss) -- (zero);
    \draw[myarrow] (zero) -- (bwd);
    \draw[myarrow] (bwd) -- (step);

    \draw[myarrow, bend right=50] (step.east) to (data.east);
    \node[mylabel] at (2.5, 0.5) {下一个epoch};

    \draw[myarrow, green!70, dashed] (2.5, 5.5) |- (eval);

    \node[mylabel, blue!70] at (2.8, 5.8) {训练阶段};
    \node[mylabel, green!70] at (5.2, 0.5) {评估阶段};
\end{tikzpicture}
\caption{完整的 PyTorch 训练循环流程。每个 epoch 包含训练阶段（数据加载、前向反向、参数更新）和周期性验证评估。}
\label{fig:training-loop}
\end{figure}

\section{模型保存与加载}

PyTorch 提供两种模型保存方式：

\begin{mycode}{模型保存与加载}
\begin{lstlisting}[language=Python]
import torch
import torch.nn as nn

model = nn.Sequential(
    nn.Linear(10, 20), nn.ReLU(), nn.Linear(20, 3))

# 方式1: 保存整个模型（不推荐，耦合度高）
torch.save(model, 'full_model.pt')
loaded_model = torch.load('full_model.pt')

# 方式2: 只保存参数（推荐）
torch.save(model.state_dict(), 'model_params.pt')

# 加载参数
new_model = nn.Sequential(
    nn.Linear(10, 20), nn.ReLU(), nn.Linear(20, 3))
new_model.load_state_dict(torch.load('model_params.pt'))

# 保存检查点（checkpoint）
checkpoint = {
    'epoch': 50,
    'model_state_dict': model.state_dict(),
    'optimizer_state_dict': optimizer.state_dict(),
    'loss': 0.0234,
    'best_acc': 0.95,
}
torch.save(checkpoint, 'checkpoint.pt')

# 恢复训练
checkpoint = torch.load('checkpoint.pt')
model.load_state_dict(checkpoint['model_state_dict'])
optimizer.load_state_dict(checkpoint['optimizer_state_dict'])
start_epoch = checkpoint['epoch']
\end{lstlisting}
\end{mycode}

\section{完整实例：MNIST手写数字识别}

\begin{mycode}{MNIST分类完整实现}
\begin{lstlisting}[language=Python]
import torch
import torch.nn as nn
import torch.optim as optim
from torch.utils.data import DataLoader
from torchvision import datasets, transforms

# 超参数
batch_size = 128
epochs = 10
lr = 0.001

# 数据准备
transform = transforms.Compose([
    transforms.ToTensor(),
    transforms.Normalize((0.1307,), (0.3081,))
])

train_dataset = datasets.MNIST(
    './data', train=True, download=True, transform=transform)
test_dataset = datasets.MNIST(
    './data', train=False, transform=transform)

train_loader = DataLoader(train_dataset, batch_size=batch_size,
                          shuffle=True)
test_loader = DataLoader(test_dataset, batch_size=batch_size,
                         shuffle=False)

# CNN模型
class CNN(nn.Module):
    def __init__(self):
        super().__init__()
        self.conv1 = nn.Conv2d(1, 32, 3, padding=1)
        self.conv2 = nn.Conv2d(32, 64, 3, padding=1)
        self.pool = nn.MaxPool2d(2, 2)
        self.fc1 = nn.Linear(64 * 7 * 7, 128)
        self.fc2 = nn.Linear(128, 10)
        self.dropout = nn.Dropout(0.25)
        self.relu = nn.ReLU()

    def forward(self, x):
        x = self.pool(self.relu(self.conv1(x)))
        x = self.pool(self.relu(self.conv2(x)))
        x = x.view(x.size(0), -1)  # 展平
        x = self.relu(self.fc1(x))
        x = self.dropout(x)
        x = self.fc2(x)
        return x

device = torch.device('cuda' if torch.cuda.is_available() else 'cpu')
model = CNN().to(device)
criterion = nn.CrossEntropyLoss()
optimizer = optim.Adam(model.parameters(), lr=lr)

# 训练
for epoch in range(epochs):
    model.train()
    running_loss = 0.0
    for batch_idx, (data, target) in enumerate(train_loader):
        data, target = data.to(device), target.to(device)

        optimizer.zero_grad()
        output = model(data)
        loss = criterion(output, target)
        loss.backward()
        optimizer.step()

        running_loss += loss.item()
        if batch_idx % 100 == 99:
            avg_loss = running_loss / 100
            print(f'Epoch {epoch+1}, Batch {batch_idx+1}, Loss: {avg_loss:.4f}')
            running_loss = 0.0

# 测试
model.eval()
correct, total = 0, 0
with torch.no_grad():
    for data, target in test_loader:
        data, target = data.to(device), target.to(device)
        output = model(data)
        pred = output.argmax(dim=1)
        correct += pred.eq(target).sum().item()
        total += data.size(0)

print(f'\n测试准确率: {100 * correct / total:.2f}%')
\end{lstlisting}
\end{mycode}

\begin{figure}[H]
\centering
\begin{tikzpicture}[
    mybox/.style={rectangle, draw=blue!60, fill=blue!5, minimum width=2.5cm, minimum height=0.8cm, font=\small, rounded corners=3pt, align=center},
    myarrow/.style={->, >=stealth, thick, blue!70}
]
    % CNN architecture
    \node[mybox] (inp) at (0, 0) {输入\\\footnotesize 28$\times$28$\times$1};
    \node[mybox] (c1) at (2.8, 0) {Conv1\\\footnotesize 28$\times$28$\times$32};
    \node[mybox] (p1) at (5.6, 0) {Pool1\\\footnotesize 14$\times$14$\times$32};
    \node[mybox] (c2) at (8.4, 0) {Conv2\\\footnotesize 14$\times$14$\times$64};
    \node[mybox] (p2) at (11.2, 0) {Pool2\\\footnotesize 7$\times$7$\times$64};
    \node[mybox, draw=red!60, fill=red!5] (flat) at (14.0, 0) {展平\\\footnotesize 3136};
    \node[mybox] (fc1) at (16.5, 0) {FC1\\\footnotesize 128};
    \node[mybox] (fc2) at (19.0, 0) {FC2\\\footnotesize 10};

    \draw[myarrow] (inp) -- (c1);
    \draw[myarrow] (c1) -- (p1);
    \draw[myarrow] (p1) -- (c2);
    \draw[myarrow] (c2) -- (p2);
    \draw[myarrow] (p2) -- (flat);
    \draw[myarrow] (flat) -- (fc1);
    \draw[myarrow] (fc1) -- (fc2);

    \node[font=\footnotesize] at (9.5, -1.2) {卷积部分\;\;$\to$\;\;全连接部分};
\end{tikzpicture}
\caption{MNIST CNN 模型结构。输入 28x28 灰度图，经过两层卷积+池化提取特征，展平后通过两层全连接层输出 10 类概率。}
\label{fig:cnn-architecture}
\end{figure}

\section{本章小结}

本章介绍了 PyTorch 的核心使用方法：
\begin{itemize}
    \item 张量是 PyTorch 的基础数据结构，支持 GPU 加速和自动微分
    \item 自动求导（Autograd）自动计算梯度，是实现反向传播的基础
    \item \texttt{nn.Module} 封装了神经网络层和参数管理
    \item PyTorch 提供了丰富的损失函数和优化器（SGD、Adam、AdamW 等）
    \item \texttt{Dataset} 和 \texttt{DataLoader} 规范了数据加载和批处理流程
    \item 完整的训练循环包括：训练模式设置、梯度清零、前向、损失计算、反向、参数更新
    \item 模型保存推荐使用 \texttt{state\_dict} 方式

    掌握这些 PyTorch 基础知识后，读者可以继续学习更高级的主题
如卷积神经网络、循环神经网络、Transformer、生成模型等。
\end{itemize}

